% titel: Blatt 4, A5: Rekursiv definierte Folge

% Teil I: Werkzeuge, kein Lösungsweg.

% Grundidee: Rekursion a_{n+1} = f(a_n); Anfangswerte; Ungleichung per
% Induktion; Monotonie über Differenz; Konvergenz; Grenzwert als
% Fixpunktgleichung L=f(L).

Eine rekursiv definierte Folge ``erzeugt sich selbst``: jedes Glied aus
dem vorigen. Um Konvergenz und Grenzwert zu finden, gehen wir in
Sch\"ritten vor: (1) erste Glieder, um ein Muster zu sehen, (2) eine
Ungleichung per Induktion, die Monotonie und Beschr\"anktheit liefert,
(3) den Grenzwert als \emph{Fixpunkt} der Rekursionsfunktion.

\begin{werkzeugkasten}[Nachrechnen der ersten Glieder]
Die Folge $a_1=1$, $a_{n+1}=\frac{1+a_n}{2+a_n}$:
\begin{equation}
  a_2=\frac2 3,\ a_3=\frac5 8,\ a_4=\frac{13}{21},\ a_5=\frac{34}{55}.
\end{equation}
Das sind Ver\"altnisse aufeinanderfolgender Fibonacci-Zahlen
$(\tfrac{1}{1},\tfrac2 3,\tfrac5 8,\tfrac{13}{21},\dots)$ --- ein
deutliches Signal, dass die Folge gegen einen festen Wert strebt
(das Verh\"altnis zweier aufeinanderfolgender Fibonacci-Zahlen strebt
gegen $\frac{\sqrt5-1}2\approx0.618$).
\end{werkzeugkasten}

\begin{werkzeugkasten}[Differenz und Monotonie]
Die Differenz zweier Folgeglieder:
\begin{equation}
  a_{n+1}-a_n=\frac{1+a_n}{2+a_n}-a_n
  =\frac{1-a_n-a_n^2}{2+a_n}.
\end{equation}
Ist $a_n^2+a_n>1$ (und $2+a_n>0$), so ist $a_{n+1}-a_n<0$: die Folge ist
streng monoton fallend. Da alle $a_n>0$, ist sie nach unten beschr\"ankt,
also konvergent (Monotonie- und Beschr\"anktheits-Kriterium).
\end{werkzeugkasten}

\begin{werkzeugkasten}[Grenzwert als Fixpunkt]
Konvergiert die rekursive Folge (d.h.\ $a_n\to L$ mit stetiger
Rekursionsfunktion), so geht die Rekursion im Grenzteil in die
\emph{Fixpunktgleichung} \"uber --- im Grenzwert wird $a_{n+1}=a_n=L$:
\begin{equation}
  L=\frac{1+L}{2+L}\ \Longleftrightarrow\ L^2+L-1=0\ \Longrightarrow\
  L=\frac{-1\pm\sqrt5}2 .
\end{equation}
Nur die positive L\"osung kommt in Frage (alle $a_n>0$ und fallend),
also $L=\frac{\sqrt5-1}2=\frac1{\Phi}$ (Kehrwert des goldenen Schnitts).
\end{werkzeugkasten}

\begin{achtung}[Darauf achten]
\begin{itemize}
  \item Beim Induktions-Beweis von $a_n^2+a_n>1$ die Rekursion so
        umbauen, dass Z\"ahler/Nenner rational zeigen, dass die
        Ungleichung bestehen bleibt; den Start $a_1$ pr\"ufen.
  \item Monotonie nicht nur ``anschaulich``: aus $a_{n+1}-a_n<0$ und
        Beschr\"anktheit (untere Schranke $0$) folgt Konvergenz.
  \item Die Fixpunktgleichung hat zwei L\"osungen; nur die positive ist
        zul\"assig ($a_n>0$).
\end{itemize}
\end{achtung}
